% GENERATED -- do not edit by hand.
The completed replicate shows how strongly the calibrated parameters depend on the measured systemic flow. Its perturbed $Q_s$ was 2.08~L/min, 39\% below the measured value; a deviation of this order is comparable to the error of pediatric Fick output computed with an assumed oxygen consumption, which has a mean bias of $+27.5\%$ and a standard deviation of 30\% against thermodilution or MRI~\cite{Ralston2023}. Recalibrating to this record fitted it as closely as the reference calibration fitted the measured record (RMSE 0.068 against 0.105), but with the systemic resistance scale changed by $+92\%$ and $E_{\mathrm{B,RV}}$ by $+141\%$. Records that differ only within the declared measurement uncertainty are therefore fitted equally well by markedly different parameter sets. This agrees with the conditioning of the sensitivity matrix of the seven-parameter set (Supplementary Section~S2) and supports calibrating such models with quantified parameter uncertainty when the clinical data are limited~\cite{Kachabi2025,Colunga2023}. The same replicate predicted RVEDV at 31.9~mL and LVEDV at 33.3~mL against 30.5 and 36.1~mL measured, where the reference calibration predicted 43.3 and 41.1~mL. The post-closure volume errors of the reference prediction are therefore of the size that pre-closure measurement uncertainty alone can produce, and a volume prediction from this record should be reported with its uncertainty, as recommended for mechanistic cardiovascular models~\cite{Colebank2024Guidelines}. The 20-replicate ensemble quantifies how often such shifts occur.
